\section{Introduction}
\label{sec:intro}

Primitive colliders are usually judged before they ever touch the simulator: by fit, count, editability, or a decomposition objective.
For robot assets, that is a premature stopping point.
A package that looks plausible against the mesh can still move a compound center of mass, destabilize contact, or erase the link structure that a Newton task needs.
The missing step is therefore not another geometric score, but an acceptance gate between a candidate proxy and an executable simulator artifact.

Recent CPD, CoACD, and V-HACD-style methods optimize geometric fidelity and proxy complexity~\cite{cpd2026,coacd2022,vhacd}.
They are valuable candidate generators, but they do not by themselves answer the operational question: should this generated package be promoted for a recorded simulation, articulation, or task setting?

We study \emph{simulation-checked} acceptance for Newton workflows.
The organizing principle is simple: treat every primitive package as a hypothesis, preserve USD provenance and link boundaries~\cite{openusd2026}, test the hypothesis with named Newton diagnostics~\cite{newton2026}, and record whether the package is accepted, rejected, or routed to fallback.
In this framing, fallback is useful evidence: it records where geometry-only candidate generation stopped being enough.

We make three contributions within a Phase~0 Newton diagnostic setting.
First, it makes the acceptance gate explicit: candidate generation, package provenance, Newton task construction, and fallback accounting are separate records.
Second, it grounds the gate in mechanism evidence: a capped bed cylinder passes in isolation but fails drop/settle as a full compound, and generated V-HACD packages for bowl, cup, and tray map into Newton yet fail named probes.
Third, it extends the same accounting to an articulated Franka smoke path, where link-framed primitives must be consumed by Newton bodies without crossing link ownership.

% Placeholder figure: compiler/checker pipeline schematic.
\begin{figure}[t]
  \centering
  \fbox{\parbox[c][1.2in][c]{0.92\linewidth}{\centering\small
    Placeholder: asset $\rightarrow$ primitive candidates $\rightarrow$ Newton diagnostics $\rightarrow$ accept / fallback}}
  \caption{Simulation-checked primitive collider workflow. Quantitative tables in the experiments section link to dated records; this schematic is not itself evidence.}
  \label{fig:pipeline}
\end{figure}

